% ============================================================
%  preamble.tex — Math for Machine Learning
%  All package imports and global settings.
%  XeLaTeX required (fontspec, unicode-math).
% ============================================================

% --- Encoding & fonts -----------------------------------------------------------
\usepackage{fontspec}
\usepackage{unicode-math}

% Main body font: Libertinus Serif (open-source; mirrors Linux Libertine)
% Fall back to Latin Modern if Libertinus is not installed.
\IfFontExistsTF{Libertinus Serif}{%
  \setmainfont{Libertinus Serif}%
}{%
  \setmainfont{Latin Modern Roman}%
}

% Sans-serif used for chapter/section headings
\IfFontExistsTF{Libertinus Sans}{%
  \setsansfont{Libertinus Sans}%
}{%
  \setsansfont{Latin Modern Sans}%
}

% Monospace used for code listings
\IfFontExistsTF{Fira Code}{%
  \setmonofont{Fira Code}[Scale=MatchLowercase]%
}{%
  \setmonofont{Latin Modern Mono}[Scale=MatchLowercase]%
}

% --- Page geometry --------------------------------------------------------------
\usepackage[
  paperwidth  = 6.125in,
  paperheight = 9.25in,
  top         = 0.875in,
  bottom      = 1in,
  inner       = 0.875in,
  outer       = 0.75in,
  includeheadfoot
]{geometry}

% --- Mathematics ----------------------------------------------------------------
\usepackage{amsmath}
\usepackage{amssymb}
\usepackage{amsthm}
\usepackage{mathtools}           % extends amsmath (pmatrix*, dcases, …)

% Per-chapter numbering of equations, figures, and tables
\usepackage{chngcntr}
\counterwithin{equation}{chapter}
\counterwithin{figure}{chapter}
\counterwithin{table}{chapter}

% Theorem-like environments
\theoremstyle{definition}
\newtheorem{definition}{Definition}[chapter]
\newtheorem{example}{Example}[chapter]

\theoremstyle{plain}
\newtheorem{theorem}{Theorem}[chapter]
\newtheorem{lemma}[theorem]{Lemma}

\theoremstyle{remark}
\newtheorem*{remark}{Remark}
\newtheorem*{note}{Note}

% --- Code listings --------------------------------------------------------------
\usepackage{listings}
\usepackage{xcolor}

\definecolor{codebg}    {HTML}{F8F8F8}
\definecolor{codeframe} {HTML}{DDDDDD}
\definecolor{keyword}   {HTML}{0000CC}
\definecolor{string}    {HTML}{006600}
\definecolor{comment}   {HTML}{888888}
\definecolor{number}    {HTML}{AA4400}

\lstdefinestyle{python}{
  language        = Python,
  backgroundcolor = \color{codebg},
  frame           = single,
  rulecolor       = \color{codeframe},
  basicstyle      = \small\ttfamily,
  keywordstyle    = \color{keyword}\bfseries,
  stringstyle     = \color{string},
  commentstyle    = \color{comment}\itshape,
  numberstyle     = \tiny\color{number},
  numbers         = left,
  numbersep       = 6pt,
  stepnumber      = 1,
  showstringspaces= false,
  breaklines      = true,
  breakatwhitespace = true,
  tabsize         = 4,
  captionpos      = b,
  aboveskip       = \medskipamount,
  belowskip       = \medskipamount,
}

\lstdefinestyle{text}{
  backgroundcolor = \color{codebg},
  frame           = single,
  rulecolor       = \color{codeframe},
  basicstyle      = \small\ttfamily,
  numbers         = none,
  showstringspaces= false,
  breaklines      = true,
  breakatwhitespace = true,
  tabsize         = 4,
  captionpos      = b,
  aboveskip       = \medskipamount,
  belowskip       = \medskipamount,
}

% Default listing style is Python
\lstset{style=python}

% --- Tables ---------------------------------------------------------------------
\usepackage{booktabs}     % \toprule, \midrule, \bottomrule
\usepackage{array}        % extended column types
\usepackage{tabularx}     % auto-width columns

% --- Figures & captions ---------------------------------------------------------
\usepackage{graphicx}
\usepackage{caption}
\usepackage{subcaption}

\captionsetup{
  font        = small,
  labelfont   = bf,
  labelsep    = period,
}

% --- Hyperlinks -----------------------------------------------------------------
\usepackage[
  colorlinks  = true,
  linkcolor   = black,
  citecolor   = black,
  urlcolor    = blue,
  pdftitle    = {Math for Machine Learning — A Software Engineer's Guide},
  pdfauthor   = {},
  pdfsubject  = {Mathematics for Machine Learning},
  pdfkeywords = {linear algebra, calculus, probability, statistics, machine learning},
  bookmarks   = true,
  bookmarksnumbered = true,
  unicode     = true,
]{hyperref}

% --- Typography -----------------------------------------------------------------
\usepackage{microtype}    % character protrusion and font expansion
\usepackage{parskip}      % inter-paragraph spacing instead of indent
\setlength{\parindent}{0pt}

% --- Block quotes ---------------------------------------------------------------
\usepackage{csquotes}     % \enquote{} and display block quotes

% --- Headers & footers ----------------------------------------------------------
\usepackage{fancyhdr}
\pagestyle{fancy}
\fancyhf{}
\fancyhead[LE]{\small\itshape\leftmark}   % even pages: chapter name
\fancyhead[RO]{\small\itshape\rightmark}  % odd pages:  section name
\fancyfoot[C]{\small\thepage}
\renewcommand{\headrulewidth}{0.4pt}

% Plain style for chapter-opening pages
\fancypagestyle{plain}{%
  \fancyhf{}%
  \fancyfoot[C]{\small\thepage}%
  \renewcommand{\headrulewidth}{0pt}%
}

% --- Misc utilities -------------------------------------------------------------
\usepackage{enumitem}     % flexible list spacing
\setlist{nosep, leftmargin=*}

\usepackage{xspace}       % smart trailing space in macros

% --- Convenience math macros ----------------------------------------------------
\newcommand{\R}{\mathbb{R}}
\newcommand{\N}{\mathbb{N}}
\newcommand{\E}{\mathbb{E}}
\newcommand{\Prob}{\mathbb{P}}
\newcommand{\Var}{\operatorname{Var}}
\newcommand{\Cov}{\operatorname{Cov}}
\newcommand{\norm}[1]{\left\lVert#1\right\rVert}
\newcommand{\abs}[1]{\left\lvert#1\right\rvert}
\newcommand{\inner}[2]{\left\langle#1,\,#2\right\rangle}
\newcommand{\T}{^{\mathsf{T}}}         % transpose: \mathbf{A}\T
\newcommand{\inv}{^{-1}}               % inverse:   \mathbf{A}\inv
